% Rationale Zahlen – bank/rationale-zahlen/ (zusammenbau v0.4, Heft)
\einheitenkopf[t5]{Rationale Zahlen}
\setcounter{aufgabe}{15}

\begin{aufgabe}{Ordnen – Prüfungsaufgaben}
\begin{teile}
\teil Gib eine Zahl an, die größer ist als $-320$. \hfill (P10 2015 OS) \\ \leerfeld
\teil Ein Tauchroboter soll höher bleiben als $-85$ m. Nenne eine mögliche Position in m. \hfill (P10 2015 OS) \\ \leerfeld[m]
\teil Ordne aufsteigend: $-\frac{2}{5}$; $1{,}3$; $-0{,}43$; $0{,}9$ \hfill (P10 2014 OS) \\ \leerfeld $<$ \leerfeld $<$ \leerfeld $<$ \leerfeld
\teil Ordne aufsteigend: $-\frac{1}{4}$; $2{,}6$; $-0{,}257$; $0{,}5$ \hfill (P10 2014 OS) \\ \leerfeld $<$ \leerfeld $<$ \leerfeld $<$ \leerfeld
\end{teile}
\end{aufgabe}

\begin{aufgabe}{Multiplizieren/Dividieren – Prüfungsaufgaben}
\begin{teile}
\teil Berechne ohne Taschenrechner den Wert von $\frac{a + b}{c}$ für $a = 15$, $b = -6$ und $c = -4$. \hfill (P10 2021 OS) \\ \leerfeld
\teil Berechne ohne Taschenrechner den Wert von $\frac{a + b}{c}$ für $a = 5$, $b = -2$ und $c = -4$. \hfill (P10 2021 OS) \\ \leerfeld
\teil Berechne den Wert von $(a + b) : c$ für $a = 3$, $b = -9$ und $c = -2$. \hfill (P10 2016 OS) \\ \leerfeld
\teil Berechne den Wert von $(a + b) : c$ für $a = 5$, $b = -19$ und $c = -7$. \hfill (P10 2016 OS) \\ \leerfeld
\teil Berechne den Wert von $3 \cdot (x - 5)$ für $x = -4$. \hfill (P10 2026 FOR) \\ \leerfeld
\teil Berechne den Wert von $6 \cdot (x - 2)$ für $x = -3$. \hfill (P10 2026 FOR) \\ \leerfeld
\teil Eine negative Zahl wird mit einer negativen Zahl multipliziert. Das Ergebnis ist … \hfill (P10 2015 OS)\\ \kreuz{immer positiv}\\ \kreuz{immer negativ}\\ \kreuz{Das kann man nicht entscheiden.}
\teil Eine negative Zahl wird durch eine positive Zahl geteilt. Das Ergebnis ist … \hfill (P10 2015 OS)\\ \kreuz{immer positiv}\\ \kreuz{immer negativ}\\ \kreuz{Das kann man nicht entscheiden.}
\end{teile}
\end{aufgabe}
